%3. Foreign aid: majority support for increased foreign aid if it is well spent
%============================================================================================================
%Reprendre/paraphraser/étendre et compléter (notamment avec mon papier NHB) la section A.1.3 de mon papier NHB
%============================================================================================================

%LB: J'ai pris quelques notes sur les baromètres que tu m'a envoyé par whattsap, si tu penses que c'est très utile de les ajouter, je peux le faire. Elle sont en dessous: 


%Eurobarometer surveys reveal broad acceptance of foreign aid across the European Union. In 2015, in the 28 EU Member States, approximately 64\% supported honouring the EU's commitment to increase aid or increasing it beyond the promised amount \citep{european_commission_2015_eurobarometer}. 
%In 2019, a survey across the 28 Member States found that 73\% of individuals considered EU should spend more on public assistance (30\%) or should continue to spend the same amount (43\%) \citep{european_commission_2019_development}. 
%In 2022, in all 27 Member States, 80\% viewed poverty reduction in developing countries as a main EU priority, and 67\% considered it a priority for their national government \citep{european_commission_2022_development}.
%In 2020, a survey of approximately across the 27 Member States found that 91\% of respondents considered EU funding for humanitarian aid important \citep{european_commission_2021_humanitarian}.

\section{Foreign aid: majority support for increased foreign aid if it is well spent}

Literature has long studied attitudes toward foreign aid in high income countries. Overall, the amount of foreign aid is overestimated by the people and a very low percentage wants, when they find it out, to cut it, in contrast with actual amounts. \citet{PIPA2001} find(s) that a multilateral initiative to eradicate world hunger by half is support by 83\% of Americans. In 20 countries, \citet{PIPA2008} indicate that a majority, 81\% on average, believe that developed countries ``have a moral responsibility to work to reduce hunger and severe poverty in poor countries''. The study shows that, minimum 75\% of respondents are willing to pay for a program aiming to eradicate world hunger by half for an estimated cost of, for instance, \$50 per year per American, in seven OECD countries. Eurobarometer data display a majority (80\%) are in favour of honouring the promise to increase aid \citep{cho_ideology_2024}.


%============================================================================================================
%- underestimation: Kaufman, Nair, Fabre et al NHB (Figs 8, S20, S23),
%============================================================================================================


\citet{kaufmann_foreign_2019} document that respondents in each of the 24 countries they study overestimate foreign aid, by a factor of seven on average. In most countries, the level of aid people would like their government to provide exceeds the level they believe is currently provided.\footnote{\citet{kaufmann_foreign_2012} provide the most reliable estimates of desired aid because, as noted by \citet{hudson_mile_2012}, earlier studies generally failed to account for misperceptions of actual aid levels.} Respondents in the top income quintile support aid levels that are 0.13 percentage points of GNI lower than those preferred by the bottom 40\%. Combining a theoretical framework with data on lobbying and observed aid allocations, while controlling for citizens' preferences in terms of desired aid, the authors conclude that actual aid falls short of desired aid because wealthier individuals exert disproportionate political influence toward their economic interests. The United States stands out as an exception. Desired aid is similar to the average across other countries, at around 3\% of GNI, but respondents believe that aid is already about twice as large as the level they would ideally support. Perceptions of current aid alone explain more than 15\% of the variation in desired aid. From one country to another, the ratio of median desired aid to actual aid ranges almost from parity to a ratio of more than 9 to 1, revealing what the authors describe as a substantial ``democratic aid deficit'' in several donor countries.

\citet{gilens_political_2001} reports that even politically informed Americans overestimate the amount of foreign aid. Providing information on actual aid levels reduces the share of respondents in favor of aid cuts by 17\%. Likewise, \citet{nair_misperceptions_2018} attribute weak support for aid in the U.S. to widespread misperceptions their rank in the global income distribution: respondents underestimate their rank by 27 percentiles on average while they overestimate the global median income by a factor of ten. Such misperceptions may explain why, in the 2000-2004 waves of the GSS, more than 60\% of Americans believed that the government was spending too much on foreign aid \citep{okten_preferences_2007}.

Using a survey experiments in the United States, the United Kingdom, France, Germany and Spain, \citet{fabre_majority_2025} show that respondents greatly overestimate current foreign aid (e.g., 4.7\% against an actual 0.4\% of public spending in the U.S.), but still prefer substantially higher aid budgets than those currently observed, allocating on average 1.8\% after receiving information on the true amount and 4.0\% without it \citet[Figure 8]{fabre_majority_2025_si} .

Respondents were also asked about their belief about foreign aid amount in their country. Most people overestimate the actual amount of foreign aid, on average, Americans beliefs actual aid is between 1.8 and 2.6\% while, in reality, the amount is 0.4\%. In France the actual level is 0.8\% while French, on average, believe this level is about 1.1\% to 1.7\% \citep[Figure S20]{fabre_majority_2025_si}.

%============================================================================================================
% - support Fabre et al NHB (Figs 6, 7, S21-S26)
%============================================================================================================

Respondents were then asked about their attitudes toward their country's evolution of foreign aid. 60\% of Americans are in favor for an increase, 43\% of them wants conditions attached to this, while the corresponding figures for France are 63\% and 51\%. These levels are close in Spain, Germany and UK \citep[Figure 6]{fabre_majority_2025_si}. In addition, respondents were asked about the government spending they would allocate to foreign aid, with a random subgroup first informed of the actual amount. Without this information, the median preferred interval was 1.8-2.6\% in the United States, France, and Germany, compared with 1.1-1.7\% in Spain and the United Kingdom. Informing respondents reduced the U.S. median to 0.3-0.5\%, while the French and German medians remained at 1.8-2.6\% \citep[Figure S20-S21]{fabre_majority_2025_si}. Respondents, for about 47\% in U.S., 59\% in Europe and even 75\% in France ``Preferred foreign aid is higher than current'' \citep[Figure S23]{fabre_majority_2025_si}. Respondents who preferred higher foreign aid were asked how the increase should be financed. ``Higher taxes on the wealthiest'' were by far the most frequently selected option, chosen by 68\% of Americans and 64\% of Europeans, and by as many as 82\% in Spain and 85\% in the United Kingdom. By contrast, relatively few, between 1\% and 8\% respondents favored cuts in education \citep[Figure S24]{fabre_majority_2025_si}. Respondents who preferred lower aid were asked how the budgetary savings should be used instead. Healthcare services spending was generally prioritized: 57\% of Europeans favored it while it represents 40\% of Americans \citep[Figure S25]{fabre_majority_2025_si}.


%============================================================================================================
% - conditions: primarily no corruption
%============================================================================================================

Respondents who are in favor of an increase in foreign aid only under certain conditions were asked: ``Under which conditions should foreign aid be increased?''. \citet{fabre_majority_2025} find that the most important condition is ensuring that ``the aid reaches people in need and money is not diverted'', supported by 77\% of Europeans and 68\% of Americans. A second widely supported condition is ``that recipient countries comply with climate targets and human rights'' (72\% in Europe and 61\% in the United States). Around half of respondents also require that ``other high-income countries also increase their foreign aid'' (51\% in Europe and 45\% in the U.S.), whereas conditioning aid on ``increased taxes on millionaires'' receives more limited support (38\% and 36\%, respectively) \citep[Figure 7]{fabre_majority_2025}.

Reviewing the literature, \citet{hudson_mile_2012} conclude that support for foreign aid is primarily motivated by concern for poverty, namely altruism. They also document skepticism about aid effectiveness. Indeed, \citet{dfid_aid_2009} and \citet{PIPA2001}, show that 47\% of British people believe that aid is wasted because of corruption, while Americans estimate that fewer than one-quarter of aid funds reach people in need and that more than half are captured by corrupted governments. Despite these beliefs, majorities continue to support foreign aid, this suggests non-utilitarian motives.

Consistent with \citet{henson_public_2010}, \citet{bauhr_corruption_2013} find that perceived corruption in recipient countries lowers support for aid. However, this relationship is mitigated by what they describe as the aid-corruption paradox: countries where corruption is major are often those with the highest need for assistance. \citet{bodenstein_who_2017} further show that right-leaning Europeans and respondents who perceive high levels of corruption are more likely to support conditioning EU development assistance on recipient countries complying with standards of democracy, human rights, and good governance.

Respondents who preferred reducing or maintaining foreign aid were asked: ``Why should foreign aid should not be increased?'' The most common reason is ``Charity begins at home: there is already a lot to do to support the [country] people in need'', selected by 63\% of both Europeans and Americans. The second most frequent answer is ``Aid is not effective as most of it is diverted'', chosen by 53\% of Europeans and 40\% of Americans \citep[Figure 7]{fabre_majority_2025}.%LB: or: \citep[Figure S25]{fabre_majority_2025_si} for the citation 


%============================================================================================================
% - determinants: higher support on the left and among rich, greater effect of ideology among the rich (Cho et al 24)
%============================================================================================================
%LB: la dernière fois, on avait parlé de (Cho, 2024) dont les résultats en intro sont les suivants: « My empirical findings indicate that financial circumstances often outweigh this normative impact of ideology, especially when respondents are explicitly directed to take their own financial situation into account when indicating the level of support they favor. My results also show that left-leaning citizens are not always more supportive of sharing resources with poor people abroad than are right-leaning voters. » (Cho, 2024, p. 542)
% --> donc je pense que je vais juste utiliser la Figure 1 page 10. Et peut être les données de l'Eurobaromètre comme tu l'avais dit (sauf si je les mobilise au début de la section). 

Using the 2002 Gallup survey and the 2006 World Values Survey, \citet{bayram_aiding_2017} and \citet{paxton_individual_2012}, consistent with the U.S. evidence of \citet{heinrich_voters_2018}, identify trust, left-wing political orientation, political interest, and female gender as the main individual-level correlates of support for higher aid.

Among left-wing voters, respondents’ economic circumstances matter, depending on whether they are lower- or higher-income. At the far left, the probability of answering totally agree'' to the statement expressing moral support for aid is about 63\% among higher-income individuals, compared with around 44\% among lower-income individuals. Conversely, lower-income respondents are more likely to choose tend to agree'' (42\%), compared with around 31\% of higher-income respondents. Therefore, the gap between richer and poorer respondents across the political spectrum is visible only on the left. Higher-income left-wing respondents express stronger support for aid, whereas lower-income left-wing respondents support it more moderately \citep{cho_ideology_2024}.